% GENERATED FILE. Edit canonical lessons, then run npm run generate:books.
% canonical-source-hash: fnv1a64:6df9a7dfe65ce6b9
% canonical-lessons: KA-C124-alugadde, KA-C124-badanekayi, KA-C124-sunthi, KA-C124-bellulli, KA-C124-jenu

\chapter{Potato, Aubergine, Ginger, Garlic, Honey}
\label{ch:ka-potato-aubergine-ginger-garlic-honey}
\hlchaptermodality{124}

\begin{chapteropening}
\textbf{By the end of this chapter:} \emph{I can say a potato, an aubergine, ginger, garlic and honey.}

Complete the last lesson of chapter 124: i can say a potato, an aubergine, ginger, garlic and honey.
\end{chapteropening}

\section[ālūgaḍḍe]{\kn{ಆಲೂಗಡ್ಡೆ} (ālūgaḍḍe) — a potato}

\label{lesson:KA-C124-alugadde}



\begin{quote}
Before the new one: say the Kannada for a coconut, then the Kannada for an onion.
\end{quote}

\subsection*{You'll want to know: \kn{ಆಲೂಗಡ್ಡೆ}}
\textbf{\kn{ಆಲೂಗಡ್ಡೆ}} — \emph{ālūgaḍḍe} — \textquotedblleft{}a potato\textquotedblright{}.

\begin{grammarlens}[title={What you've built}]
One more word for this chapter.
\end{grammarlens}

\subsection*{Your turn}
\begin{itemize}
\raggedright
  \item \emph{Say it:} \emph{ālūgaḍḍe}
  \item \emph{Say it:} \emph{ālūgaḍḍe}, once more
  \item \emph{Say it:} say \emph{ālūgaḍḍe}
  \item \emph{Recall:} say the Kannada for a coconut, then the Kannada for an onion, then say \emph{ālūgaḍḍe} again
\end{itemize}

\begin{tcolorbox}[breakable,colback=teal!4,colframe=teal!35!black,title={Before you move on}]
What does \kn{ಆಲೂಗಡ್ಡೆ} mean? (\textquotedblleft{}A potato\textquotedblright{}.) Say it once more.
\end{tcolorbox}
\section[badanekāyi]{\kn{ಬದನೆಕಾಯಿ} (badanekāyi) — an aubergine}

\label{lesson:KA-C124-badanekayi}



\begin{quote}
Before the new one: say the Kannada for an onion, then the Kannada for a potato.
\end{quote}

\subsection*{You'll want to know: \kn{ಬದನೆಕಾಯಿ}}
\textbf{\kn{ಬದನೆಕಾಯಿ}} — \emph{badanekāyi} — \textquotedblleft{}an aubergine\textquotedblright{}.

\begin{grammarlens}[title={What you've built}]
One more word for this chapter.
\end{grammarlens}

\subsection*{Your turn}
\begin{itemize}
\raggedright
  \item \emph{Say it:} \emph{badanekāyi}
  \item \emph{Say it:} \emph{badanekāyi}, once more
  \item \emph{Say it:} say \emph{badanekāyi}
  \item \emph{Recall:} say the Kannada for an onion, then the Kannada for a potato, then say \emph{badanekāyi} again
\end{itemize}

\begin{tcolorbox}[breakable,colback=teal!4,colframe=teal!35!black,title={Before you move on}]
What does \kn{ಬದನೆಕಾಯಿ} mean? (\textquotedblleft{}An aubergine\textquotedblright{}.) Say it once more.
\end{tcolorbox}
\section[śuṇṭhi]{\kn{ಶುಂಠಿ} (śuṇṭhi) — ginger}

\label{lesson:KA-C124-sunthi}



\begin{quote}
Before the new one: say the Kannada for a potato, then the Kannada for an aubergine.
\end{quote}

\subsection*{You'll want to know: \kn{ಶುಂಠಿ}}
\textbf{\kn{ಶುಂಠಿ}} — \emph{śuṇṭhi} — \textquotedblleft{}ginger\textquotedblright{}.

\begin{grammarlens}[title={What you've built}]
One more word for this chapter.
\end{grammarlens}

\subsection*{Your turn}
\begin{itemize}
\raggedright
  \item \emph{Say it:} \emph{śuṇṭhi}
  \item \emph{Say it:} \emph{śuṇṭhi}, once more
  \item \emph{Say it:} say \emph{śuṇṭhi}
  \item \emph{Recall:} say the Kannada for a potato, then the Kannada for an aubergine, then say \emph{śuṇṭhi} again
\end{itemize}

\begin{tcolorbox}[breakable,colback=teal!4,colframe=teal!35!black,title={Before you move on}]
What does \kn{ಶುಂಠಿ} mean? (\textquotedblleft{}Ginger\textquotedblright{}.) Say it once more.
\end{tcolorbox}
\section[beḷḷuḷḷi]{\kn{ಬೆಳ್ಳುಳ್ಳಿ} (beḷḷuḷḷi) — garlic}

\label{lesson:KA-C124-bellulli}



\begin{quote}
Before the new one: say the Kannada for an aubergine, then the Kannada for ginger.
\end{quote}

\subsection*{You'll want to know: \kn{ಬೆಳ್ಳುಳ್ಳಿ}}
\textbf{\kn{ಬೆಳ್ಳುಳ್ಳಿ}} — \emph{beḷḷuḷḷi} — \textquotedblleft{}garlic\textquotedblright{}.

\begin{grammarlens}[title={What you've built}]
One more word for this chapter.
\end{grammarlens}

\subsection*{Your turn}
\begin{itemize}
\raggedright
  \item \emph{Say it:} \emph{beḷḷuḷḷi}
  \item \emph{Say it:} \emph{beḷḷuḷḷi}, once more
  \item \emph{Say it:} say \emph{beḷḷuḷḷi}
  \item \emph{Recall:} say the Kannada for an aubergine, then the Kannada for ginger, then say \emph{beḷḷuḷḷi} again
\end{itemize}

\begin{tcolorbox}[breakable,colback=teal!4,colframe=teal!35!black,title={Before you move on}]
What does \kn{ಬೆಳ್ಳುಳ್ಳಿ} mean? (\textquotedblleft{}Garlic\textquotedblright{}.) Say it once more.
\end{tcolorbox}
\section[jēnu]{\kn{ಜೇನು} (jēnu) — honey}

\label{lesson:KA-C124-jenu}



\begin{quote}
Before the new one: say the Kannada for ginger, then the Kannada for garlic.
\end{quote}

\subsection*{You'll want to know: \kn{ಜೇನು}}
\textbf{\kn{ಜೇನು}} — \emph{jēnu} — \textquotedblleft{}honey\textquotedblright{}.

\begin{grammarlens}[title={What you've built}]
One more word for this chapter.
\end{grammarlens}

\subsection*{Your turn}
\begin{itemize}
\raggedright
  \item \emph{Say it:} \emph{jēnu}
  \item \emph{Say it:} \emph{jēnu}, once more
  \item \emph{Say it:} say \emph{jēnu}
  \item \emph{Recall:} say the Kannada for ginger, then the Kannada for garlic, then say \emph{jēnu} again
\end{itemize}

\begin{tcolorbox}[breakable,colback=teal!4,colframe=teal!35!black,title={Before you move on}]
What does \kn{ಜೇನು} mean? (\textquotedblleft{}Honey\textquotedblright{}.) Say it once more.
\end{tcolorbox}
